\chapter{Arhitectura rețelei}
\section{Experimente}
În această secțiune sunt descrise diferite tipuri de rețele utilizate și arhitecturile utilizate, comparând rezultatele obținute și motivând
alegerea arhitecturii finale.
\subsection{CNN}
Rețelele  neuronale convoluționale (CNNs) sunt rețele neuronale utilizate, de obicei, pentru procesarea imaginilor. 
Prima aplicație ce are la bază rețelele convoluționale moderne \cite{cnns-first} a demonstrat faptul că o astfel de rețea 
poate recunoaște cifrele scrise de mână din setul de date MNIST cu o rată de eroare de doar 0.7\% (Boosted LeNet-4), un rezultat
semnificativ mai bun decât alte rețele ce nu aveau la bază straturi convoluționale și, de asemenea, cu un număr mai mic de parametrii
ce pot fi antrenați.

Ideea de bază a acestui tip de rețea este emularea modului în care oamenii procesează imaginile, prin 
identificarea unor structuri simple care agregate, definesc structuri mai complexe iar în final obiecte. Acest proces se bazează pe 
două elemente fundamentale, convoluția si agregarea (pooling).

\subsubsection{Convoluția}
Un strat convoluțional are la bază unul sau mai multe nuclee (kernels). Fiecare nucleu este o matrice, bidimensională sau
tridimensională, în funcție de numărul de canale al imaginii. Se calculează un produs scalar dintre nucleu și o submatricea 
imaginii de aceeași dimensiune cu cea a nucleului, rezultatul produsului fiind salvat in celula corespunzătoarea celulei 
care se află în centrul submatricii selectate. Procesul este ilustrat în figura \ref{fig:f15} 
\footnote{\href{https://github.com/musikalkemist/DeepLearningForAudioWithPython/blob/master/15\-\%20How\%20does\%20a\%20convolutional\%20neural\%20network\%20work/slides/How\%20does\%20a\%20convolutional\%20\%20neural\%20network\%20work.pdf}{How does a convolutional neural network work}}. 
Pentru a calcula produsul 
scalar și pentru elementele de pe marginea matricii, în cazul în care nu se dorește reducerea dimensiunii, matricea va fi bordată.

\begin{figure}[h]
    \caption{Produsul scalar dintre submatrice și nuclee}
    \includegraphics[width=\textwidth]{kernels.png}
    \centering
    \label{fig:f15}
\end{figure}

Nucleul se deplasează cu un anumit număr de poziții (strides) pe verticală și pe orizontală. Acest factor poate modifica, din nou, 
dimensiunea rezultatelor. Fiecare nucleu din stratul convoluțional va genera ca rezultat o matrice bidimensională. 
Prin antrenare, nucleele acținează ca filtre ce vor fi modificate pentru a identifica structurile de baza ale imaginii, 
linii și curbe în primele 
straturi convoluționale și structuri mai complexe la adâncimi mai mari.

După un strat convolutional, se aplică o operație de agregare (pooling).

\subsubsection{Pooling}

Pooling-ul este o operație ce are ca scop reducerea dimensiunii reprezentării. Similar cu procesul de convoluție, 
se stabilește o dimensiune de pooling și un stride. Diferența este faptul că, operația nu depinde de parametri ce pot 
fi antrenați, este doar un calcul asupra elementelor dintr-o submatrice cu aceeași dimensiune ca dimensiune de pooling 
stabilită. Această operație poate fi selectarea elementului maxim din submatrice (Max Pooling), calcularea mediei, 
calcularea mediei ponderate în funcție de distanța de la centrul matricii, etc.

Procesul de Max Pooling pentru o 
dimensiune de (2, 2) și un stride de 2 este ilustrat în figura \ref{fig:f16}

\begin{figure}[h]
    \caption{Max Pooling pentru dimensiunea (2, 2) și un stride de 2}
    \includegraphics[width=0.5\textwidth]{pooling.png}
    \centering
    \label{fig:f16}
\end{figure}

\subsubsection{Batch Normalization}
Batch Normalization este un proces prin care input-urile straturilor rețelei sunt normalizate, rezultatul fiind o antrenare mai 
stabilă și rapidă a rețelei. 

Normalizarea se obține calculând media batch-ului,
$\mu_{B} = \frac{1}{m}\sum^{m}_{i=1}x_{i}$, și deviația standard, $\sigma^{2}_{B} = \frac{1}{m}\sum^{m}_{i=1}(x_{i} - \mu_{B})^{2}$, 
unde $m$ este numărul de elemente din batch. Se calculează $\hat x_{i}=\frac{x_{i}-\mu_{B}}{\sqrt{\sigma^{2}_{B}+\epsilon}}$, iar pe baza
acestor valori se calculează noile valori normalizate, $y_{i}=\gamma\hat x_{i} + \beta$, unde $\gamma$ și $\beta$ sunt paramentrii ce
pot fi antrenați. \cite{batch-norm}

\subsubsection{Arhitectura CNN implementată}

Având în vedere că spectrogramele Mel și coeficienții MFCC pot fi considerate reprezentări similare cu o 
imagine cu un singur canal (grayscale), ce mențin aspectele de structură prezente în imagini, rețelele convoluționale 
sunt populare nu doar în procesarea imaginilor, ci și în procesarea audio.
Arhitectura rețelei convoluționale implementate are 3 straturi convoluționale cu 32 de nuclee de dimensiunea (3, 3). 
Fiecare strat convoluțional este urmat de un strat de Max Pooling cu dimensiunile (3, 3), (3, 3), respectiv (2, 2) apoi de 
un strat de Batch Normalization. Output-urile sunt aplatizate și introduse ca input într-un strat dens conectat cu 64 de neuroni, urmat de 
stratul de output. 

\subsubsection{Dropout}
Dropout-ul este un procedeu prin care, în timpul antrenării, un procent din 
neuronii rețelei dintr-un anumit layer sunt ignorați, alături de conexiunile 
cu ceilalți neuroni. Acești neuroni sunt
selectați random la fiecare iterație. Rezultatul este prevenirea fenomenului
 de overfitting deoarece devine imposibil ca rețeaua să devină dependentă
de un subset al neuronilor, atunci când aceștia pot deveni inactivi. Astfel,
este incurajată generalizarea. \cite{overfitting}

Din acest motiv, între straturile dens conectate am utilizat un strat de Dropout de 0.3.
\subsubsection{ReLU}
Funcțiile de activare pentru straturile convoluționale și primul strat dens sunt ReLU, $f(x) = max(0, x)$ deoarece
previne problema de dispariție a gradientului (vanishing gradient) și are o complexitate computațională scăzută.
\subsubsection{Softmax}
Ultimul strat este unul dens conectat cu un număr de neuroni egal cu numărul distincte de genuri și activare Softmax.
Fiecare valoare a output-ului va reprezenta un procent, sugestiv pentru gradul în care
un anumit gen se regăsește întru-un segment. 
\subsubsection{Adam}
Optimizatorul utilizat a fost Adam cu o rată de învățare de 0.005 deoarece s-a demonstrat că are rezultate mai bune în
practică decât stochastic gradient descent \cite{adam}.

\pagebreak
\subsection{LSTM}
\subsubsection{Rețele neuronale recurente}
Rețelele neuronale recurtente sunt un tip de rețele neuronale în care conexiunile dintre
noduri formează un graf orientat în care informația circulă în sensul axei temporale, prin
existența unor stări și inputuri ascunse. \cite{rnn}
Structura generală a unei rețele recurente este ilustrată în figura \ref{fig:f17} \footnote{\href{https://stanford.edu/~shervine/teaching/cs-230/cheatsheet-recurrent-neural-networks}{Recurrent Neural Networks cheatsheet}}.

\begin{figure}[h]
    \caption{Arhitectura tradițională unei rețele neuronale recurente}
    \includegraphics[width=0.8\textwidth]{rnn.png}
    \centering
    \label{fig:f17}
\end{figure}

$a^{<t>}=g_{1}(W_{aa}a^{<t-1>}+W_{ax}x^{<t>}+b_{a})$ și $y^{<t>} = g_{2}(W_{ya}a^{<t>} + b_{y})$, unde $W_{aa},\ W_{ax}$ și
$W_{ya}$ sunt ponderile pentru conexiuni (weights) iar $b_{a}$ și $b_{y}$ sunt valorile bias. Atât ponderile, cât și valorile
bias sunt parametrii ce vor fi antrenați.

Rețelele neuronale recurente sunt des utilizate în problemele de procesarea limbajului natural, 
de speech recognition și chiar de generarea de segmente muzicale, probleme în care ordinea informațiilor
are o mare importanță, deoarece iau în calcul aspectul
temporal al datelor.   

O problemă a acestor tip de rețele este dificultatea accesării informațiilor care 
se află un număr mai mare de pași în trecut deoarece gradientul poate crește sau
descrește exponențial odată cu numărul mare de straturi care crește odată cu lungimea
datelor de intrare.

\subsubsection{LSTM}
Rețelele Long Short Term Memory (LSTM) sunt un subtip de rețele neuronale recurente
care rezolvă problema învățării dependențelor de lungă durată, prin adăugarea
unei stări menite să rețină informații pentru perioade mai lungi de timp.

Structura rețelei este similară cu cea a unei rețele recurente tradiționale, diferența
fiind la nivelul celulei care, acum nu mai conține doar un strat, ci patru straturi care
interacționează unele cu altele \ref{fig:f18} \footnote{\label{note1}\url{http://colah.github.io/posts/2015-08-Understanding-LSTMs/}}. 

\begin{figure}[h]
    \caption{Arhitectura LSTM}
    \includegraphics[width=0.8\textwidth]{lstm.png}
    \centering
    \label{fig:f18}
\end{figure}

Zona roșie a figurii reprezintă elementele ce țin de memoria de lungă durată
a rețelei. Se observă faptul că asupra acestei stări se execută doar două operații.
Prima operație este legată de uitarea(eliminarea) informațiilor, iar a doua operație
este asociată cu adăugarea de noi informații.

Prima operație este o înmulțire bazată pe elemente (produs Hadamard) dintre 
starea de scurtă durată curentă si output-ul stratului sigmoid, $f_{t} = \sigma(W_{f}\cdot[h_{t-1}, x_{t}] + b_{f})$, unde $[h_{t-1}, x_{t}]$ reprezintă
concatenarea output-ului stratului anterior cu input-ul stării curente.
Datorită proprităților funcției de activare sigmoidă, valorile rezulate aparțin intervalului $[0, 1]$, unde $1$
semnifică faptul că informațiila va fi păstrată, iar pentru $0$ va fi complet eliminată \cite{lstm}.

A doua operație asupra stării ce ține de memoria de lungă durată este ilustrată în figura \ref{fig:f19} \ref{note1}.

\begin{figure}[h]
    \caption{Actualizarea stării celulei}
    \includegraphics[width=0.7\textwidth]{lstm1.png}
    \centering
    \label{fig:f19}
\end{figure}

Similar cu prima operație efectuată output-ul stratului sigmoid, $i_{t}$, va
acționa ca un filtru asupra valorilor $\widetilde{C}_{t}$ ce vor fi adunate la 
starea $C_{t}$ a celulei. Astfel, se adaugă stării doar informațiile noi ce sunt
considerate relevante de rețea.

Zona verde din figura \ref{fig:f18} reprezintă output-ul celulei. Acest output va fi 
construit prin aplicarea unei funcții de activare (tangenta hiperbolica) asupra stării celulei
și aplicarea unui filtru asupra rezultatului prin înmulțirea cu output-ul stratului sigmoid.

\subsubsection{Arhitectura LSTM implementată}
Având în vedere contextul temportal al datelor extrase din segmentele audio, 
a fost construit un model bazat pe LSTM pentru a vedea dacă o procesare ce ține
cont de contextul temporal reusește să aibă performanțe mai bune decât o
rețea convoluțională, ținând cont de datele problemei abordate.

Modelul construit este format din doua straturi LSTM cu 64 de unități, urmate de un strat
dens conectat cu 64 de neuroni, un strat Dropout de 0.3 și un strat dens conectat de
output cu activare Softmax și un număr de neuroni egal cu numărul de genuri distinte 
din setul de date.

\subsection{CNN-LSTM}

Combinarea rețelelor convoluționale cu rețele recurente este des utilizată în practică pentru 
recunoașterea activităților din videoclipuri sau din secvențe de imagini, dar și pentru recunoașterea
sunetelor sau a genurilor muzicale. 

Rețelele CNN-LSTM combină elementele de procesare a datelor bidimensionale sau
tridimensionale ale rețelelor convoluționale și aspectele temporale surprinse de
rețelele LSTM. Straturile convoluționale reduc dimensiunea inputurilor și extrag, în 
același timp caracteristicile esențiale ale datelor de intrare. Caracteristicile
extrase de rețeaua CNN sunt apoi introduse într-o rețea LSTM care procesează evoluția
caracteristicilor, generând o predicție.

\subsubsection{Arhitectura CNN-LSTM implementată}
Rețeaua CNN are 4 straturi convoluționale cu 64 de nuclee de dimensiunea (3, 3). Fiecare strat
convoluțional este urmat de un strat de Max Pooling cu dimensiunile (2, 2), (3, 3), (4, 4), (4, 4), 
și un strat de Batch Normalization. Output-ul rețelei CNN este aplatizat și introdus ca input în
rețeaua LSTM. 

Rețeaua LSTM are aceeasi arhitectură cu cea implementată în capitolul anterior. În continuare 
sunt prezentate comparațiile între rețelele implementate.
\subsection{Comparații}
\subsection{Selecția genurilor}
\section{Concluzii}